% =====================================================================
\chapter{导论}
% =====================================================================

\section{问题的由来}

以下三个问题指向同一结构。它们最初以工程直觉的形式提出；本节给出可讨论、可证明、
可证伪的形式。

\begin{problem}[底层确定性的可推广性]\label{prob:low}
软件底层的运行过程具有高度的一致性：ABI 与调用约定、函数调用与栈帧、堆与对象生命
周期、控制流与数据流、同步原语与线程调度。这些机制在不同软件之间几乎相同，其一致性
来自共同的计算模型、硬件架构、语言运行时与操作系统接口。问题是：这种确定性是否可以
推广到高层？大部分正常的、偏置不多的大型复杂软件系统，是否也会在抽象、结构与架构上
收敛到相同或相似的形态？
\end{problem}

\begin{problem}[规模的强制性]\label{prob:scale}
当软件代码量或协作复杂度跨过某一量级，系统是否会\emph{必然}开始针对同类问题进行
设计与优化，从而在内部长出一套统一的基础设施？Linux 内核会定义内存如何管理、如何
分配、进程与线程如何定义与调度，并为上层提供运行时接口；浏览器、游戏引擎、异步
运行时、分布式存储与关系数据库引擎也会考虑同一类问题，只是所处的层次各不相同。
代码量一旦上去，这些系统是否都会在代码中体现出"将功能统一、接口统一、管理规划、
运行时结构"这类思考？
\end{problem}

\begin{problem}[智能的理解与设计]\label{prob:mind}
软件是人类与 AI 编写的，但软件本身并不容易理解。大型复杂软件几乎不可能从微观开始
顺序阅读理解。那么应当从什么抽象层级理解？如何由理解进入开发，而不陷入代码海洋？
智能不能陷入无穷的底层信息量，而须保留对底层基本原理的理解
（记为 \Lk{0}）与上层的 \Lk{1}、\Lk{2}、\Lk{3}……，即被压缩过的信息。
反过来，究竟应当如何设计一个大型复杂软件系统，使目的得以达成？
\end{problem}

第三个问题有一个尖锐的当代变体，它把前两个问题从"如何理解"推向"如何保证"：

\begin{problem}[无人审阅条件下的正确性]\label{prob:trust}
如果 AI 使编码成本趋近于零，人不再逐行审阅代码，那么如何\emph{机械地}、\emph{非自指地}
保证正确性？即：如何保证人提出的初始目标被实现，且实现过程中没有额外的、未被授权的
副作用？一个具体的例子是：原本手写一次调用耗时 1\,ns，AI 生成的实现耗时 1\,$\mu$s，
功能完全一致，目的达成，却产生了性能上的副作用。
\end{problem}

\section{为什么这些问题尚未被统一处理}

上述问题中的每一个都在既有文献中有局部对应：架构收敛对应模式语言与复杂系统理论，
规模的强制性对应信息隐藏、近可分解性与模块化理论，理解与设计对应认知模型与多视图
建模，正确性对应程序逻辑、精化演算与验证编译。但把它们作为一个整体处理时，会出现
三处结构性缺口。

\paragraph{缺口一：收敛命题的强度未被区分。}
"大型软件架构会收敛"这句话至少包含三个强度差异极大的命题：局部架构模式的反复出现
（有大量工程经验支持）；在适当抽象下呈现高于随机基线的可测量结构相似性（很可能成立，
值得系统验证）；几乎所有正常运行的大型软件最终收敛到同一种高层架构（证据不足，且
可能过强）。多数讨论把三者混为一谈，因而无法被检验。

\paragraph{缺口二：必然性主张缺少承重结构。}
"规模强制产生统一基础设施"是一个必然性主张，但通常以经验观察的形式陈述，因而无法
回答：必然性从何而来？阈值由什么决定？为什么有些同样巨大的系统没有长出内核？本文
的答案是：必然性来自\emph{操作语义的完备性}，阈值来自\emph{摊销不等式}，而"没有长出
内核"的系统是把内核\emph{外借}给了宿主。

\paragraph{缺口三：理解、设计与保证是同一结构的三次出现。}
理解被当作认知问题，设计被当作工程问题，保证被当作逻辑问题，三者分别在不同的文献
传统中处理。三者是同一个对象的不同读取方向：层级就是\emph{保证义务的分解
单元}，理解是恢复这套分解，设计是规定这套分解，而保证是在这套分解上逐层施加约束。

\section{核心论题}

本文以四条论题组织全部内容。每条论题后面括号内标注了它在本文中的形式化位置与
强度等级：\textbf{[P]} 表示有完整证明，\textbf{[A]} 表示有严格论证但依赖明确假设，
\textbf{[H]} 表示作为可证伪假设提出。

\begin{principle}[压缩论题]\label{th:compression}\textbf{[A]}
一个复杂系统可被有限智能理解，当且仅当它是可压缩的。理解被定义为\emph{任务约束下的
有损压缩}：在保持任务相关预测能力的约束下最小化描述长度。由此，下潜的最优深度是
率失真拉格朗日量的驻点条件，即边际预测增益为零处。
\end{principle}

\begin{principle}[必然性论题]\label{th:necessity}\textbf{[A]}
任何具有并发存活实体与受争用有限资源的系统，若其操作语义良定义，则必然已定义六项
职能：身份与生命周期、分配与回收、调度、命名与寻址、通信与同步、故障隔离与恢复。
这六项职能\emph{可以被自建，也可以被外借}给操作系统、语言运行时或框架，但不可以在
语义上缺失，缺失即意味着存在未定义行为。
\end{principle}

\begin{principle}[分离论题]\label{th:separation}\textbf{[A]}
生成可以自由、概率、开源、不可信；准入必须确定、机械、可穷尽检视。形式层（规格
$\to$ 实现 $\to$ 二进制 $\to$ 执行）的保证来自极小可信内核的机械判定；意图层
（意图 $\to$ 规格）的保证来自对抗性共识，且这一层不存在机械解。两层由制品绑定
缝合：只有被验证过的那个二进制被允许运行，且环境假设在其上被检查或监控。
\end{principle}

\begin{principle}[设计论题]\label{th:design}\textbf{[A]}
架构就是人为制造的尺度分离。自然界的层级由尺度分离与重整化\emph{免费}获得；软件的
层级不会自然涌现，必须由接口、封装与不变量\emph{付费}制造，并且随时会漏。因此设计
是正向压缩，理解是逆向压缩，两者是同一结构的两个读取方向。架构的可测量目标函数是
一次典型变更所需的\emph{读集}期望 $\mathbb{E}[\rho]$。
\end{principle}

四条论题及其形式化结果合称\term{压缩—保证框架}。名称取自主张的两端：\emph{压缩}
是理解与设计的机制，\emph{保证}是可信性的锚定方式。框架的核心断言是：
\emph{一个复杂系统能被有界主体理解，当且仅当它可被压缩；而它可被压缩，当且仅当它
存在一套尺度分离，使每一层都拥有一个被显式声明的无视集。}

\section{认识论纪律}

本文处理的材料跨越工程直觉、复杂系统理论、信息论、数理逻辑与社会认识论，越界风险
很高。为控制风险，全文遵守以下纪律。

\begin{enumerate}[label=(\roman*)]
\item \term{强度标注}。每个主张标注其为定理、命题、猜想、假设或启发式。拒绝把启发式
      写成定理。
\item \term{不作主张清单}。第 \ref{ch:objections} 章显式列出本文\emph{不}主张的内容，
      包括被本文拒绝的若干"权威定理"。
\item \term{假设显式化}。凡结论依赖环境假设，必须在陈述中写明；凡保证有边界，必须在
      边界处标注。
\item \term{可证伪性}。第 \ref{ch:program} 章把核心命题降格为十二条可证伪假设，并给出
      操作化、对照组、证伪条件与风险。
\item \term{不与结论同构的验证}。凡提出"生成不可信、验证需独立"的地方，该原则同样
      适用于本文自身：本文的核心承重命题必须可被独立检验，且不得依赖任何未证明猜想。
\end{enumerate}

第 (v) 条约束本文自身。若干修辞性主张被刻意从承重结构中卸下：\emph{不}用 Gödel
不完备定理或 Rice 定理去证明"意图与规格之间存在原则性裂缝"
（理由见第 \ref{ch:intent} 节 \ref{sec:intent-gap}），\emph{不}用 $\mathrm{P}\neq
\mathrm{NP}$ 的未决性去支撑"验证比生成便宜"（理由见第 \ref{ch:intent} 节
\ref{sec:asymmetry}），也\emph{不}把 Simon 的近可分解性当作普遍定律（见
第 \ref{ch:convergence} 章与第 \ref{ch:objections} 章）。这几条由卸载而非加固获得
强度：拆掉未经证明的引用之后，核心命题不再依赖任何猜想。

\section{主要贡献}

\paragraph{形式化贡献。}
\begin{enumerate}[label=(C\arabic*)]
\item 以抽象映射的\emph{核}刻画层级的无视集，使"这一层可以不关心什么"从修辞变为
      可计算对象（定义 \ref{def:oblivious}）。
\item 以 Galois 连接刻画层间语义保持，并给出层级可靠性的条件式（定义
      \ref{def:galois}、命题 \ref{prop:soundness}）。
\item 以"接口是模块行为的\emph{充分统计量}"给出 Simon 近可分解性的可检验判据，并由
      此构造\emph{泄漏对}诊断程序（定义 \ref{def:neardecomp}、命题
      \ref{prop:leakpair}）。
\item 以集成成本不等式 $\Theta(nm)\to\mathcal{O}(n+m)$ 给出内部基础设施化的相变阈值，
      并说明阈值由"独立特性数 $\times$ 共享状态数"而非代码行数决定（定理
      \ref{th:integration}）。
\item 以操作语义完备性论证内核六职能的\emph{条件必然性}，并由此导出"自建 vs 外借"
      的判定准则（定理 \ref{th:six}、推论 \ref{cor:buildadopt}）。
\item 以抽象映射的像空间维度给出过约束/欠约束判据，从而把"底层唯一、高层多解"
      形式化为解集几何（命题 \ref{prop:overconstrain}）。
\item 以读集期望 $\mathbb{E}[\rho]$ 给出架构质量的可测量目标函数，并证明不变量集中
      化最小化读集（定义 \ref{def:readset}、定理 \ref{th:centralise}）。
\item 以描述长度差给出技术债的精确定义（定义 \ref{def:debt}）。
\item 以 Hoeffding 界给出共识层可靠性的指数上界，并证明当存在相关错误模式 $q$ 时该界
      被 $q$ 支配、增加节点数无益，从而确立"多样性而非数量"为共识设计的核心变量
      （定理 \ref{th:condorcet}、命题 \ref{th:correlated}）。
\item 给出完整的保证条件式，把功能、资源、副作用、系统组合与目标表达五类正确性统一
      为一条可追溯的条件蕴含（定义 \ref{def:assurance}）。
\end{enumerate}

\paragraph{方法论贡献。}
\begin{enumerate}[label=(M\arabic*)]
\item 一份七步"下行协议"，把理解大型系统从顺序阅读改造为分层建模与主动辨识。
\item 一组可执行的"缝探测器"清单，用于在陌生代码库中定位真实模块边界。
\item 一份七步"正向压缩"设计流程，把设计目标从形容词改造为可优化目标函数。
\item 十二条可证伪假设及实际上验设计。
\item 一份把社会认识论与证明携带代码耦合起来的交叉研究方向，即文献地图上最空白的
      可占据位置。
\end{enumerate}

\section{全文结构}

第 \ref{ch:framework} 章建立统一的形式框架：层级、抽象映射、无视集、契约、轨迹、
保证与信任。第 \ref{ch:convergence} 章处理问题 \ref{prob:low}，把收敛拆成三个强度
不同的命题并给出强度定律。第 \ref{ch:infrastructure} 章处理问题 \ref{prob:scale}，
证明六职能的条件必然性并给出相变阈值。第 \ref{ch:comprehension} 章处理问题
\ref{prob:mind} 的前半，给出逆向压缩的理解理论。第 \ref{ch:design} 章处理问题
\ref{prob:mind} 的后半，给出正向压缩的设计理论。第 \ref{ch:intent} 章处理问题
\ref{prob:trust}，给出从意图到可验证执行的七段链与双层信任结构。第
\ref{ch:levels} 章把层级从认知辅助升级为验证义务的分解单元。第 \ref{ch:program}
章给出可证伪的研究纲领。第 \ref{ch:literature} 章给出文献地图与定位。第
\ref{ch:objections} 章给出反对意见与限度。第 \ref{ch:conclusion} 章收束。

\begin{remark}[阅读指南]
只关心架构收敛的读者可读第 \ref{ch:framework}、\ref{ch:convergence}、
\ref{ch:infrastructure} 章；只关心理解与设计的读者可读第 \ref{ch:comprehension}、
\ref{ch:design} 章；只关心正确性与验证的读者可读第 \ref{ch:intent}、
\ref{ch:levels} 章；关心可检验性的读者应直接读第 \ref{ch:program} 章。
\end{remark}

\section{术语约定}

本文使用若干在中英文语境中含义漂移的术语，此处一次性约定，完整术语表见附录
\ref{app:glossary}。

\begin{longtable}{L{3.1cm}L{11.4cm}}
\toprule
\textbf{术语} & \textbf{本文含义} \\
\midrule
\endhead
层级 / level & 一组（表示空间、行为集合、观察映射、环境假设、成本模型、不变量），
  相对研究对象与当前任务定义，不是代码的物理目录。 \\
无视集 / oblivious set & 某一层在原理上永远无法区分的差异集合；等于该层抽象映射的核。 \\
尺度分离 & 上下层之间能够各自成立有效描述、且下层的实现细节不在上层语义中出现。 \\
内部基础设施化 & 系统把重复出现的管理问题提升为统一机制、并让上层依赖这些机制的过程。 \\
内核六职能 & 身份与生命周期、分配与回收、调度、命名与寻址、通信与同步、故障隔离与恢复。 \\
读集 / read-set & 完成一次典型变更所必须阅读的模块集合。 \\
保证 / assurance & 在明确前提与明确范围内，某性质成立的、可追溯的证据关系。 \\
机械保证 & 由一个极小、确定、可穷尽检视的判定器给出的保证。 \\
共识保证 & 由多个独立、对抗性参与方给出的、统计学意义上可靠的判断。 \\
制品绑定 & 被验证的工件与实际执行的工件之间的同一性关系。 \\
意图鸿沟 & 自然语言（或更一般的）意图与其形式化投影之间的失真。 \\
\bottomrule
\end{longtable}

一个术语上的重要提醒：本文所说的"层"是\emph{认知尺度}，而非物理层级。同一个组件
在不同研究任务中可以处于不同层：一个异步运行时对普通应用开发者是基础设施，对其自身
开发者则是需要深入修改的复杂系统。层级是相对于研究对象与当前任务定义的（见
第 \ref{ch:levels} 章）。
